\documentclass[a4paper,12pt]{article}

% Pacotes para codificação e idioma
\usepackage[utf8]{inputenc}
\usepackage[T1]{fontenc}
\usepackage[brazil]{babel}

% Pacotes matemáticos essenciais
\usepackage{amsmath}
\usepackage{amssymb}
\usepackage{amsfonts}
\usepackage{mathtools}

% Pacote para formatação de margens
\usepackage{geometry}
\geometry{a4paper, left=3cm, right=2cm, top=3cm, bottom=2cm}

% Título e Autor
\title{\textbf{Formulação Matemática: Drenagem Profunda em Córregos com Aduelas em Meio Não Saturado}}
\author{Modelagem Numérica}
\date{\today}

\begin{document}

\maketitle

\section{Introdução}
Este documento descreve o modelo matemático para a simulação de fluxo em meios porosos variadamente saturados, aplicado ao contexto de drenagem profunda utilizando aduelas (box culverts). A física do problema é governada pela conservação de massa acoplada à Lei de Darcy generalizada.

\section{Equação Governante}
O fluxo de água na zona vadosa (solo não saturado) e na zona saturada é descrito pela \textbf{Equação de Richards}. Adota-se aqui a formulação mista (carga de pressão e conteúdo de umidade) para garantir melhor conservação de massa numérica.

A equação governante em um domínio $\Omega \subset \mathbb{R}^d$ ($d=2$ ou $3$) é dada por:

\begin{equation}
    \frac{\partial \theta(h)}{\partial t} = \nabla \cdot \left[ \mathbf{K}(h) \cdot (\nabla h + \mathbf{e}_z) \right] + S
    \label{eq:richards}
\end{equation}

\noindent onde:
\begin{itemize}
    \item $h$: Carga de pressão [L] (variável primária).
    \item $\theta(h)$: Conteúdo volumétrico de umidade [-].
    \item $t$: Tempo [T].
    \item $\mathbf{K}(h)$: Tensor de condutividade hidráulica [L/T].
    \item $\mathbf{e}_z$: Vetor unitário na direção vertical (gravidade), $\mathbf{e}_z = (0, \dots, 1)^T$.
    \item $S$: Termo de fonte/sumidouro volumétrico [1/T].
\end{itemize}

A carga hidráulica total $H$ é definida como $H = h + z$, onde $z$ é a cota de elevação.

\section{Relações Constitutivas (Modelo de Solo)}
Para fechar o sistema de equações, as propriedades hidráulicas do solo ($\theta$ e $K$) são descritas pelas curvas características de retenção de água e condutividade. Utiliza-se o modelo de \textbf{Van Genuchten-Mualem}:

\subsection{Saturação Efetiva ($S_e$)}
\begin{equation}
    S_e(h) = 
    \begin{cases} 
        \left[ 1 + |\alpha h|^n \right]^{-m} & \text{se } h < 0 \\
        1 & \text{se } h \ge 0
    \end{cases}
\end{equation}
onde $m = 1 - 1/n$, e $\alpha, n$ são parâmetros empíricos de forma do solo.

\subsection{Conteúdo de Umidade ($\theta$)}
\begin{equation}
    \theta(h) = \theta_r + (\theta_s - \theta_r) S_e(h)
\end{equation}
onde $\theta_s$ é o conteúdo de umidade saturado (porosidade) e $\theta_r$ é o conteúdo residual.

\subsection{Função de Condutividade Hidráulica ($K$)}
\begin{equation}
    K(h) = K_s \cdot S_e^{1/2} \left[ 1 - \left( 1 - S_e^{1/m} \right)^m \right]^2
\end{equation}
onde $K_s$ é a condutividade hidráulica do solo saturado.

\section{Condições de Contorno}
A modelagem correta da aduela requer condições de contorno específicas para representar a interface solo-estrutura e a drenagem (seepage face).

\subsection{Fronteira do Córrego ($\Gamma_{c}$)}
Assume-se uma condição de Dirichlet baseada no nível da água do córrego $z_w$:
\begin{equation}
    h(\mathbf{x}, t) = \max(0, z_w - z) \quad \text{em } \Gamma_{c}
\end{equation}

\subsection{Fronteira da Aduela/Dreno ($\Gamma_{d}$)}
A parede da aduela atua como uma \textit{superfície de infiltração}. A água só pode sair do solo para a aduela se a pressão for positiva (saturada); caso contrário, a parede é impermeável. Isso configura um problema de condição de contorno complementar (tipo Signorini):

\begin{equation}
    \begin{cases} 
        h \le 0 \\
        Q_n \ge 0 \quad (\text{Fluxo saindo do domínio}) \\
        h \cdot Q_n = 0
    \end{cases} \quad \text{em } \Gamma_{d}
\end{equation}

Numericamente, isso é implementado verificando a cada iteração:
\begin{itemize}
    \item Se o nó está saturado ($h \ge 0$): Impõe-se $h = 0$ (Pressão Atmosférica).
    \item Se o nó está não-saturado ($h < 0$): Impõe-se $\frac{\partial h}{\partial n} = 0$ (Fluxo Nulo).
\end{itemize}

\subsection{Fronteiras Impermeáveis ($\Gamma_{imp}$)}
Para a base do aquífero ou divisores de águas laterais:
\begin{equation}
    \mathbf{q} \cdot \mathbf{n} = -\mathbf{K}(h) \nabla H \cdot \mathbf{n} = 0 \quad \text{em } \Gamma_{imp}
\end{equation}

\section{Discretização Temporal}
Devido à alta não-linearidade de $C(h) = \partial \theta / \partial h$, recomenda-se o uso de um esquema de Euler Implícito (Backward Euler) com linearização via método de Picard ou Newton-Raphson:

\begin{equation}
    C(h^{k}) \frac{h^{k+1} - h^{t}}{\Delta t} - \nabla \cdot \left[ \mathbf{K}(h^k) (\nabla h^{k+1} + \mathbf{e}_z) \right] = S
\end{equation}
Onde $k$ denota a iteração dentro do passo de tempo atual.

\end{document}
